\documentclass[11pt,letterpaper]{article}
\usepackage[margin=1in,headheight=15pt]{geometry}
\usepackage{fontspec}
\setmainfont{Latin Modern Roman}
\setsansfont{Latin Modern Sans}
\setmonofont{DejaVu Sans Mono}[Scale=0.78]
\usepackage{amsmath,amssymb,mathtools,amsthm}
\usepackage{microtype}
\usepackage{xcolor,booktabs,tabularx,longtable,array,enumitem}
\usepackage{fvextra}
\usepackage{fancyhdr,titlesec}
\usepackage{hyperref,bookmark}
\definecolor{ink}{HTML}{18364A}
\definecolor{accent}{HTML}{176B78}
\definecolor{muted}{HTML}{52616B}
\hypersetup{colorlinks=true,linkcolor=accent,urlcolor=accent,citecolor=accent,
 pdftitle={A Constructive Lean Contender Beyond Contender 5},
 pdfauthor={ChatGPT}}
\titleformat{\section}{\Large\sffamily\bfseries\color{ink}}{\thesection}{0.7em}{}
\titleformat{\subsection}{\large\sffamily\bfseries\color{ink}}{\thesubsection}{0.7em}{}
\setlist{itemsep=3pt,topsep=5pt}
\setlength{\emergencystretch}{2em}
\pagestyle{fancy}
\fancyhf{}
\fancyhead[L]{\small\sffamily Constructive Lean Contender}
\fancyhead[R]{\small\sffamily Research and checked implementation}
\fancyfoot[C]{\small\thepage}
\renewcommand{\headrulewidth}{0.3pt}
\newtheorem{theorem}{Theorem}[section]
\newtheorem{lemma}[theorem]{Lemma}
\newtheorem{proposition}[theorem]{Proposition}
\newtheorem{corollary}[theorem]{Corollary}
\theoremstyle{definition}
\newtheorem{definition}[theorem]{Definition}
\newtheorem{remark}[theorem]{Remark}
\newcommand{\N}{\mathbb N}
\newcommand{\llbracket}{\mathopen{[\![}}
\newcommand{\rrbracket}{\mathclose{]\!]}}
\newcommand{\Cfive}{\ensuremath{C_5}}
\newcommand{\Cnew}{\ensuremath{N_\star}}
\newcommand{\eval}{\operatorname{eval}}
\newcommand{\code}[1]{\texttt{#1}}
\newcommand{\artifact}[1]{\texttt{\detokenize{#1}}}
\newcommand{\verified}{\textsf{Lean checked}}
\begin{document}
\begin{titlepage}
\vspace*{1.0cm}
{\sffamily\color{accent}\Large NAME THE BIGGEST NUMBER\par}
\vspace{0.7cm}
{\sffamily\bfseries\color{ink}\Huge A Constructive Lean Contender\\[0.25cm]Beyond Contender 5\par}
\vspace{0.6cm}
{\Large Total interpretation, well-founded trees, and feedback through power towers\par}
\vspace{1.0cm}
{\large A reproducible Lean 4 contest entry\par}
{\large 12 October 2026\par}
\vspace{1.1cm}
\noindent\rule{\textwidth}{0.5pt}
\begin{abstract}
We adapt Cody Roux's constructive largest-number contest from Coq to Lean 4 and
supply a concrete, executable-in-principle natural number, \Cnew, together with
a kernel-checked proof that it strictly exceeds the repository's Contender 5.
The baseline is the original total interpretation of raw simply typed
lambda-calculus syntax with natural-number recursion, maximized over all
terms of depth at most 42. Its Lean translation preserves malformed terms,
computational casts, reversed variable indices, and the original winner
selection rule.

Our construction combines a cumulative envelope of this language with
structural recursion on countably branching well-founded trees. An exact
grafting identity turns substitution of trees into replacement of the base
growth function. Higher function spaces and explicitly specified closure
constructors supply the tree payload; score-dependent iteration supplies a
quantitative guarantee independent of any ordinal interpretation. With
\(T(0)=1\), \(T(h+1)=2^{T(h)}\), and
\(\Theta(s,0)=s\), \(\Theta(s,k+1)=T(\Theta(s,k)+1)\), the checked result is
\[
 \boxed{\quad \Cfive<\Cnew,
 \qquad \Theta(\Cfive+1,\Cfive+1)\leq\Cnew.\quad}
\]
All 32 audited declarations have empty transitive axiom lists. Fresh cumulative
checking of all custom definitions and all comparison proofs fits the contest's
15-second and one-minute budgets on the recorded machine. The construction is
an explicit submission, with no claim of optimality among all admissible Lean
programs. Exact ordinal calibration of the Veblen-inspired syntax is not used;
an explicit counterexample explains why such calibration cannot be assumed.
\end{abstract}
\vfill
{\small\color{muted}Pinned implementation: Lean 4.19.0.\quad
Upstream snapshot: \texttt{d56dce013e62}.\\
The archive contains the article, all Lean sources, an independent small-input
model, reproduction scripts, source provenance, and validation records.}
\end{titlepage}

\tableofcontents
\newpage

\section{The proposed entry and the claim it establishes}

The submitted number is the closed Lean definition
\begin{Verbatim}[frame=single,fontsize=\small]
namespace BigNumber

def contender_6 : Nat := Amplifier.candidate base

end BigNumber
\end{Verbatim}
This short declaration names a fully specified construction; all dependencies
are included and counted in the validation budget. The main proof is named
\code{BigNumber.contender\_5\_lt\_contender\_6}. A stronger theorem,
\code{BigNumber.contender\_6\_tower\_bound}, establishes the boxed inequality
in the abstract. Neither theorem evaluates \Cfive\ or \Cnew.

The lower bound has a direct reading. Start at \(\Cfive+1\). Construct a tower
of twos whose height is one greater than that starting value. Use the resulting
integer, plus one, as the height of the next tower. Repeat this operation
\(\Cfive+1\) times. The resulting integer is a proved lower bound for the
submitted number. The actual definition also contains substantial tree and
higher-type structure that this estimate does not attempt to measure.

There are three distinct levels of assurance:
\begin{enumerate}
\item The mathematical definitions and proofs below explain the construction
      without requiring the reader to understand Lean tactics.
\item The supplied Lean files check the exact translated baseline, tree
      algebra, and numerical comparison with no axiom dependencies.
\item A source audit connects the translated baseline to the Coq source.
      This last connection is a reviewed translation, supported by independent
      small-input tests; it is not a formal Coq-to-Lean compiler-correctness
      theorem.
\end{enumerate}

The source repository was read at commit
\begin{center}
\code{d56dce013e6264292533efdc5e1de71863c6fb86}.
\end{center}
The computational definitions in \code{Contender.v}, including the designated
fifth contender, are the benchmark throughout this article~\cite{contender}.

\section{A precise Lean adaptation of the rules}

The original rules require constructive numbers, formal comparison proofs,
no axioms, a 15-second definition budget, and a one-minute proof budget. They
also forbid defining a new entry by referring to the previous closed
contender and ask for a substantial margin~\cite{rules}. We retain that intent
and make the Lean-specific checks explicit.

\begin{center}
\begin{tabularx}{\textwidth}{@{}p{0.23\textwidth}X@{}}
\toprule
Requirement & Lean interpretation used here\\
\midrule
A natural number & A closed, safe, computable definition of type \code{Nat},
using ordinary inductive recursors.\\[3pt]
Definition budget & At most 15 seconds cumulatively to check all custom
computational dependencies from source. The official standard library may be
precompiled.\\[3pt]
Comparison budget & At most 60 seconds cumulatively for the custom proofs;
we additionally include the axiom and dependency audit in the measured proof
budget.\\[3pt]
No axioms & Empty transitive axiom lists for the entry and the advertised
comparison theorems. No \code{sorry}, declared axioms, classical choice,
user-written unsafe definitions, or noncomputable data.\\[3pt]
Substantive new entry & No transitive dependence of the submitted number on
the previous closed contender. General language definitions and total
interpreters remain reusable mathematical infrastructure.\\[3pt]
Reproducibility & Pin Lean 4.19.0, disclose the source snapshot, preserve
original licensing, and include commands and recorded validation.\\
\bottomrule
\end{tabularx}
\end{center}

The allotted times concern \emph{checking the definition and proof}, not
computing the decimal expansion and not the time spent designing the entry.
The user supplied no separate deadline for development. We therefore used the
repository's explicit checking budgets.

Lean's logical framework includes proof irrelevance as a kernel conversion
rule. Additional principles such as propositional extensionality, quotient
soundness, and classical choice are visible to the axiom audit. The official
manual explains the distinction between pure computation, proof-related
extensions, and noncomputable choice~\cite{leanaxioms}. This submission stays
within the strict empty-dependency interpretation. Importing a library that
contains an unused axiom does not itself make a declaration depend on it.

A few convenience lemmas in Lean's library carry \code{propext} transitively.
The supplied proof files replace the particular maximum and multiplication
lemmas that would have introduced that dependency with direct inductive
proofs. These replacements are short; there is no need to weaken the contest's
no-axioms rule.

\section{What Contender 5 actually denotes}
\label{sec:baseline}

\subsection{Raw syntax and depth}

The source defines simple types by
\[
 A,B ::= \mathsf{Nat}\mid A\to B
\]
and raw terms by
\[
 t,u ::= \mathsf{var}(i)\mid\lambda^{A,B}t\mid t\,u
          \mid 0\mid S\mid\mathsf{rec}_A.
\]
The lambda constructor carries both a domain and a codomain annotation.
Raw syntax is deliberately not restricted by a typing judgment. Every raw
term receives a total, default-valued interpretation.

Type depth is
\[
 d(\mathsf{Nat})=1,\qquad d(A\to B)=1+\max(d(A),d(B)).
\]
The term-depth conventions are important:
\begin{align*}
 d(\mathsf{var}(i))&=i+2,\\
 d(\lambda^{A,B}t)&=1+\max(d(A),d(B),d(t)),\\
 d(tu)&=1+\max(d(t),d(u)),\\
 d(0)&=d(S)=1,\qquad d(\mathsf{rec}_A)=1+d(A).
\end{align*}
Let \(\mathcal R_n\) be the exact list produced by the source's
\code{termsUpTo n}. The original file proves that this list contains exactly
the raw terms of depth at most \(n\). Our Lean port reproduces its recursion,
product nesting, concatenation order, and duplicate behavior literally.

\subsection{A total semantic interpreter}

The interpretation of a simple type is the corresponding natural-number or
function type in the host theory. A value is a dependent pair
\(\langle A,v\rangle\) with \(v\in\llbracket A\rrbracket\).
The error value is zero at natural type and a constant error function at an
arrow type. The interpreter is structural recursion on the raw syntax.

The following details prevent accidental replacement of the benchmark by an
easier, different number:
\begin{itemize}
\item Variable zero denotes the final environment element. With an environment
      of length \(\ell\), variable \(i<\ell\) selects position
      \(\ell-1-i\). A newly bound value is prepended. Thus indices are levels
      counting from the outermost binder, not ordinary nearest-binder indices.
\item A lambda evaluates its body in the extended environment and casts the
      result to the annotated codomain. Both annotations affect raw evaluation.
\item An application whose head has natural type returns the natural error
      value. An arrow-valued head casts its argument into the domain before
      applying the function.
\item Equal arrow types use computational, recursively eta-expanded casts.
      Incompatible types produce the appropriate default value. The definition
      does not transport values through an opaque equality proof.
\item The recursor at \(A\) has type
      \(A\to(\N\to A\to A)\to\N\to A\). Its step receives the predecessor
      index and the previous value.
\item After evaluating a closed raw term, a natural-valued result is read as
      its integer; a function-valued result is read as zero.
\end{itemize}

These rules cover malformed terms, not just the well-typed fragment usually
called G\"odel's system~T. No normalization theorem for untyped lambda calculus
is being assumed. The total host-language definition itself supplies the
computation.

\subsection{The selected maximum and the numerical envelope}

The source traverses \(\mathcal R_n\), keeping a current best term. It replaces
the winner only when the new value is \emph{strictly} greater, so the first
winner is retained on ties. Define the numerical envelope
\[
 E(n)=\operatorname{foldmax}_{t\in\mathcal R_n}\eval(t),
 \qquad\text{with initial maximum }0.
\]
The literal port of the source's closed entry is
\[
 \Cfive=\eval(\operatorname{maxBy}(\eval,0,0,\mathcal R_{42})).
\]

\begin{lemma}[Selected maximum agrees with the envelope]
\label{lem:envelope}
For every \(n\), the value selected by the original winner-selection
algorithm is \(E(n)\). In particular, \(\Cfive=E(42)\).
\end{lemma}
\begin{proof}
Use the stronger invariant that the current winning term has value exactly
the current maximum \(m\). For an empty list, the selected term has value
\(m\), which is the fold's initial accumulator. For a head value \(v\), either
\(m<v\), in which case the new winner has value \(v=\max(m,v)\), or
\(v\le m\), in which case the existing winner has value
\(m=\max(m,v)\). Induction on the remaining list completes the proof.
The initial winning term is zero and has value zero. The Lean theorem
\code{eval\_maxBy} proves the invariant, and
\code{contender\_5\_eq\_envelope} specializes it.
\end{proof}

Enumeration alone grows rapidly. If \(a_n\) is the number of types and
\(u_n\) the number of raw terms, counting the constructors gives
\begin{align*}
 a_0=u_0&=0, & a_{n+1}&=1+a_n^2,\\
 &&u_{n+1}&=n+a_n^2u_n+u_n^2+2+a_n.
\end{align*}
The two implementations agree on the following tractable prefix:
\begin{center}
\begin{tabular}{rrrr}
\toprule
Depth & Types & Raw terms & Envelope\\
\midrule
0&0&0&0\\
1&1&2&0\\
2&2&10&1\\
3&5&146&2\\
4&26&24,976&3\\
\bottomrule
\end{tabular}
\end{center}
For example, \(u_{n+1}\ge u_n^2\) and \(u_1=2\), so
\(u_{42}\ge2^{2^{41}}\). This is a lower bound on the size of the list,
not a claimed estimate of its largest evaluated number. It explains why
expanding the depth-42 list has no role in checking the submission.

\section{An independent inflationary base}
\label{sec:base}

The new number uses the generic interpreter and envelope, but has no dependency
on the closed definition \code{contender\_5}. We define
\[
 P(0)=E(0),\qquad P(n+1)=P(n)+E(n+1),\qquad b(n)=P(n)+n+1.
\]
Equivalently, \(P(n)=\sum_{j=0}^{n}E(j)\). The finite sum provides a convenient
monotone envelope without requiring a separate monotonicity proof for the
original syntax enumerator.

\begin{proposition}[Base domination]
\label{prop:base}
For every \(n\), \(n+1\le b(n)\). For \(k\le n\),
\(E(k)<b(n)\). Consequently \(\Cfive<b(64)\).
\end{proposition}
\begin{proof}
Each envelope is a natural number. Hence \(P(n)\le P(n+1)\), and
\(E(k)\le P(k)\). Induction along \(k\le n\) yields
\(E(k)\le P(n)<P(n)+n+1=b(n)\). Also \(n+1\le b(n)\)
because \(P(n)\ge0\). Set \(k=42,n=64\) and use
Lemma~\ref{lem:envelope}.
\end{proof}

The implementation is eleven lines of computational code plus direct order
proofs. The numeral 64 is a fixed input and starting higher-type level. The
later feedback chooses much larger finite levels from previously computed
scores; 64 is not a truncation of those later stages.

\section{A constructive amplifier on well-founded trees}
\label{amp:construction}

The amplification mechanism is independent of the benchmark enumerator.  Its
input is an explicitly defined function $b:\mathbb N\to\mathbb N$ satisfying
$b(n)\geq n+1$.  Its output is a closed natural number $N_b$.  The construction
has two complementary features.  Countably branching well-founded trees and
operators on higher function spaces provide a substantial supply of recursive
structure.  A separate feedback mechanism supplies a quantitative growth
certificate that does not require assigning ordinal names to those trees.
The result is an exact definition together with an independently checkable
lower bound.

Throughout this section, iteration means ordinary finite iteration:
\[
 f^{[0]}(x)=x,
 \qquad f^{[r+1]}(x)=f^{[r]}(f(x)).
\]
This convention exposes the first application of $f$.  It defines the same
iterate as the familiar recurrence $f^{[r+1]}(x)=f(f^{[r]}(x))$.
The brackets distinguish iteration from exponentiation.

\subsection{Trees, evaluation, and substitution}
\label{amp:trees}

Let $\mathcal T$ be the inductively generated type with constructors
\[
 0_{\mathcal T},\qquad
 S(a)\quad(a\in\mathcal T),\qquad
 L(f)\quad(f:\mathbb N\to\mathcal T).
\]
The third constructor packages a countable family of immediate subtrees.
The word ``inductively'' is essential: these are well-founded trees, and their
eliminator permits structural recursion on a tree, including recursive calls
to any member $f(r)$ of a family appearing at a limit node.  They are not
coinductive trees, and a cyclic equation does not introduce a new inhabitant.
We use $0_{\mathcal T}$ rather than the numeral $0$ when its type might be
ambiguous.

For a base $b:\mathbb N\to\mathbb N$, define
\begin{align}
 F_b(0_{\mathcal T};n)&=b(n),\label{amp:eq:fastzero}\\
 F_b(S(a);n)&=\bigl(F_b(a;\mathord\cdot)\bigr)^{[n+1]}(n),
 \label{amp:eq:fastsucc}\\
 F_b(L(f);n)&=F_b(f(n);n).
 \label{amp:eq:fastlimit}
\end{align}
Thus a successor node turns the function supplied by its child into a finite
iterator.  A limit node selects its $n$th child.  The argument $n$ is allowed
to change inside a successor iteration; consequently the evaluator genuinely
uses the functions associated with subtrees, rather than merely counting
nodes along one fixed path.

For $a,t\in\mathcal T$, define the substitution $a\star t$ by replacing every
zero leaf of $t$ with $a$:
\begin{align*}
 a\star0_{\mathcal T}&=a,\\
 a\star S(t)&=S(a\star t),\\
 a\star L(f)&=L(r\mapsto a\star f(r)).
\end{align*}
The orientation matters: the first argument is inserted into the second.
The Lean implementation calls this operation \texttt{graft}.

The evaluator returns a function of a natural argument at each tree.  Its
successor equation can therefore reuse the child function with increasingly
large arguments, although the recursive call defining that function is
still to a smaller tree.  This distinction is useful when reading the
termination argument: the natural arguments to child functions need not
decrease.  Structural recursion decreases the tree argument, while the
separate iteration operation decreases its finite iteration counter.
The construction never asks the evaluator to recurse on a tree that it has
just generated from itself.

\begin{lemma}[Inflationarity and the base bound]
\label{amp:lem:base}
If $b(n)\geq n$ for every $n$, then for every $a\in\mathcal T$ and every $n$,
\[
 n\leq b(n)\leq F_b(a;n).
\]
If instead $b(n)\geq n+1$ for every $n$, then
$F_b(a;n)\geq n+1$ for every $a,n$.
No monotonicity hypothesis on $b$ is required.
\end{lemma}

\begin{proof}
First observe that $g(x)\geq x$ implies $g^{[r]}(x)\geq x$ for all $r,x$;
this follows by induction on $r$.  Now use structural induction on $a$.
The zero case is equality.  In the successor case, put
$g(x)=F_b(a;x)$.  The induction hypothesis gives $g(x)\geq b(x)\geq x$.
Therefore
\[
 F_b(S(a);n)=g^{[n]}(g(n))\geq g(n)\geq b(n).
\]
At a limit node, use the induction hypothesis for the selected child $f(n)$.
This proves the weak statement.  The strict statement follows by first
weakening $b(n)\geq n+1$ to $b(n)\geq n$ and then retaining the stronger first
inequality in $F_b(a;n)\geq b(n)\geq n+1$.
\end{proof}

The absence of a monotonicity assumption is more than a simplification of
the statement.  An inflationary function may take a smaller value at a
larger argument; nevertheless its own orbit is nondecreasing, because each
successive orbit value is at least its immediate predecessor.  The proof
uses precisely this orbit property.  It does not replace $b(x)$ by $b(y)$
merely because $x\leq y$.  This makes the interface applicable to bases
assembled from finite searches whose dependence on their search parameter
has not separately been proved monotone.

\begin{theorem}[Substitution law]
\label{amp:thm:graft}
For every base $b$, without an inflationarity assumption, every $a,t$, and
every $n$,
\[
 F_b(a\star t;n)=F_{F_b(a;\mathord\cdot)}(t;n).
\]
\end{theorem}

\begin{proof}
Proceed by structural induction on $t$.  At a zero leaf both sides are
$F_b(a;n)$.  For a successor node, the induction hypothesis identifies the
two child functions at every natural argument.  Pointwise equal functions
have pointwise equal finite iterates: prove this auxiliary assertion by
induction on the iteration count, replacing the first application and then
the remaining iterate.  Applying it with count $n+1$ gives the successor
case.  At a limit node, both evaluators select the child indexed by $n$, and
the induction hypothesis for that child applies directly.  The proof uses
pointwise equalities throughout.  In particular, the Lean proof does not
invoke function extensionality.
\end{proof}

\begin{proposition}
\label{amp:cor:graftlower}
If $b(n)\geq n$ everywhere, then
\[
 F_b(a\star t;n)\geq F_b(a;n)
 \qquad(a,t\in\mathcal T).
\]
\end{proposition}

\begin{proof}
Lemma~\ref{amp:lem:base} says that $F_b(a;\mathord\cdot)$ is an inflationary
base.  Apply the same lemma to that base and to $t$, and then use
Theorem~\ref{amp:thm:graft}.
\end{proof}

The substitution law is the main algebraic safeguard in the construction.
It permits an elaborate tree to be attached to an existing tree while
preserving a proved evaluation bound.  The later operators need not satisfy
an unproved comparison principle for ordinal notations.

More specifically, $t$ acts as a program transforming the initial unary
function supplied at its zero leaves.  Substituting $a$ into those leaves
replaces that initial function by the already constructed function
$F_b(a;\mathord\cdot)$.  The theorem identifies these two operations exactly,
not merely up to a lower bound.  The tree $t$ may depend on $a$, on an
evaluation of $a$, or on other previously defined data.  Once those data
have been supplied, the same substitution theorem applies to the resulting
pair of trees.  That flexibility is what permits the later feedback
construction without weakening its proof obligation.

\begin{proposition}[The zero argument and successor gain]
\label{amp:prop:gain}
For every base $b$ and tree $a$,
\[
 F_b(a;0)=b(0).
\]
If $b(n)\geq n+1$ everywhere, then
\[
 F_b(S(a);n)\geq F_b(a;n)+n.
\]
\end{proposition}

\begin{proof}
The first assertion follows by tree induction.  A successor performs exactly
one child application at argument zero, and a limit selects its zeroth
child, so both cases reduce to the corresponding induction hypothesis.
For the second assertion, an induction on $r$ gives
$g^{[r]}(x)\geq x+r$ whenever $g(x)\geq x+1$ for all $x$.
Apply this to $g=F_b(a;\mathord\cdot)$, which has the required property by
Lemma~\ref{amp:lem:base}, and calculate
\[
 F_b(S(a);n)=g^{[n]}(g(n))\geq g(n)+n.
\]
The gain vanishes at $n=0$, exactly as the first assertion requires.
\end{proof}

For a small concrete calibration, take $b(n)=n+1$ and finite successor trees
$a_0=0_{\mathcal T}$, $a_{r+1}=S(a_r)$.  Directly from the definition,
\[
 F_b(a_0;n)=n+1,\qquad F_b(a_1;n)=2n+1,
\]
and
\[
 F_b(a_2;n)=2^{n+1}(n+1)-1.
\]
Indeed, the $r$fold iterate of $x\mapsto2x+1$ at $n$ is
$2^r(n+1)-1$, as an induction on $r$ verifies.  Substituting $r=n+1$ gives
the displayed formula.  These elementary calculations concern the
successor-only fragment; they do not supply an ordinal interpretation of
the general limit constructor.
The next tree iterates the last function $n+1$ times.  These identities
explain the basic acceleration before countable branching or higher
function spaces enter the definition.  They also show why the evaluation
argument is part of the construction, rather than an incidental input.

\subsection{Explicit tree operators}
\label{amp:operators}

We now specify the entire additional tree payload.  The following definitions
are syntactic operations on $\mathcal T$.  Their resemblance to familiar
ordinal constructions guides the design, but no identification with a
Veblen ordinal or an ordinal collapsing function is assumed in any proof.

For a tree transformation $f:\mathcal T\to\mathcal T$, set
\[
 O_f(a)=L\bigl(r\mapsto f^{[r+1]}(a)\bigr).
\]
This \emph{orbit-limit constructor} packages successively longer finite
orbits.  The name does not assert that it produces a fixed point of an
arbitrary $f$.

Define a tree transformation $W$ by
\begin{align*}
 W(0_{\mathcal T})&=S(0_{\mathcal T}),\\
 W(S(a))&=L\bigl(r\mapsto(x\mapsto W(a)\star x)^{[r]}(0_{\mathcal T})\bigr),\\
 W(L(f))&=L(r\mapsto W(f(r))).
\end{align*}
The implementation names this transformation \texttt{omegaPower}.  In
particular, its successor clause iterates substitution into a zero tree;
it is not an appeal to an externally supplied ordinal exponentiation
operation.

For a supplied transformation $c:\mathcal T\to\mathcal T$, define
\begin{align*}
 E_c(0_{\mathcal T})&=c(0_{\mathcal T}),\\
 E_c(S(a))&=c(S(E_c(a))),\\
 E_c(L(f))&=L(r\mapsto E_c(f(r))).
\end{align*}
This is an enumeration constructor generated by $c$.  For a countable family
$\boldsymbol f=(f_0,f_1,\ldots)$ of tree transformations, also define finite
prefix applications and repeated sweeps:
\begin{align*}
 P_{\boldsymbol f}(0,a)&=a,&
 P_{\boldsymbol f}(r+1,a)&=f_r(P_{\boldsymbol f}(r,a)),\\
 Q_{\boldsymbol f}(0,a)&=a,&
 Q_{\boldsymbol f}(r+1,a)&=
 P_{\boldsymbol f}(r+1,Q_{\boldsymbol f}(r,a)).
\end{align*}
The common-limit constructor is
\[
 C_{\boldsymbol f}(a)=L(r\mapsto Q_{\boldsymbol f}(r+1,a)).
\]
Repeated sweeps deliberately revisit earlier members of the family.  This
choice is more structured than applying each $f_r$ once and then discarding
it.  However, our correctness argument only needs the explicit recurrence;
it does not require a fixed-point theorem for $C_{\boldsymbol f}$.

Define a two-argument transformation $V$ by recursion on its first tree:
\begin{align*}
 V(0_{\mathcal T},a)&=W(a),\\
 V(S(t),a)&=E_{O_{V(t,\mathord\cdot)}}(a),\\
 V(L(f),a)&=
 E_{C_{(V(f(r),\mathord\cdot))_{r\in\mathbb N}}}(a).
\end{align*}
Here $O_{V(t,\mathord\cdot)}$ is the transformation
$x\mapsto O_{V(t,\mathord\cdot)}(x)$, and similarly for the common-limit
constructor in the last line.  The recursive calls $V(t,\mathord\cdot)$ and
$V(f(r),\mathord\cdot)$ involve immediate subtrees of the first argument.
Finally, set
\[
 \gamma=O_{a\mapsto V(a,0_{\mathcal T})}(0_{\mathcal T}).
\]
The symbol $\gamma$ denotes this particular tree.  It is not a claim that
the tree has been proved to represent $\Gamma_0$.

These definitions produce one explicit transformation rather than leaving
a black-box ``large ordinal'' in the candidate.  Every tree used below is
built from the three constructors and the recurrences just displayed.
Their role in the size certificate is equally explicit: a resulting tree
is placed in the second argument of substitution, where
Proposition~\ref{amp:cor:graftlower} applies regardless of its further behavior.

The dependency order also explains totality of this part of the definition.
Substitution is available before $W$ is defined.  The orbit, enumeration,
prefix, and sweep operations are available before $V$ is defined.  In the
successor clause of $V$, the transformation associated with the smaller
first tree has already been supplied by structural recursion; enumeration
then acts on the second tree using that transformation as a parameter.
The limit clause receives the analogous family of transformations for the
children of its first tree.  There is no recursive call to $V$ on a newly
constructed first argument in these clauses.

\subsection{A dynamically chosen tower of function spaces}
\label{amp:levels}

The next layer iterates operators at successively higher function types.
Define
\[
 \mathcal L_0=\mathcal T,
 \qquad \mathcal L_{k+1}=\mathcal L_k\to\mathcal L_k.
\]
Thus an element of $\mathcal L_{k+2}$ transforms endomorphisms of
$\mathcal L_k$.  Each index is a natural number, and the family itself is
defined by natural-number recursion.  No operation indexed by an infinite
natural number is being introduced.

Successor and countable branching extend pointwise through these spaces.
For $x\in\mathcal L_k$ and $f:\mathbb N\to\mathcal L_k$, define
\begin{align*}
 s_0(x)&=S(x),&
 s_{k+1}(g)(x)&=s_k(g(x)),\\
 \ell_0(f)&=L(f),&
 \ell_{k+1}(f)(x)&=\ell_k(r\mapsto f(r)(x)).
\end{align*}
If $g\in\mathcal L_{k+1}$, put
\[
 \Omega_k(g)(x)=
 \ell_k\bigl(r\mapsto g^{[r+1]}(s_k(x))\bigr).
\]
The successor $s_k(x)$ is inserted before the finite orbit.  This detail is
part of the exact definition and is retained at every function level.
The operator $\Omega_k$ belongs to $\mathcal L_{k+2}$.

Let
\[
 B_0=W,\qquad B_{k+1}=\Omega_k.
\]
In both cases $B_k\in\mathcal L_{k+1}$.  A lowering operation converts the
result back into a transformation of trees:
\[
 D_0(f)=f,\qquad
 D_{k+1}(f)=D_k(f(B_k)).
\]
The types match at each step: if $f\in\mathcal L_{k+2}$, then
$f(B_k)\in\mathcal L_{k+1}$.  Define the seed family
\[
 \Sigma_k=D_k(B_k)(\gamma).
\]
This family is defined for every natural $k$.  Later, $k$ will itself be an
earlier evaluation score plus $64$.  The dependent type family therefore
provides a compact description of a finite function-space tower whose
height is determined by the construction, rather than fixed by a long
list of manually written arrow types.

The first three seeds make the difference between levels concrete:
\begin{align*}
 \Sigma_0&=W(\gamma),\\
 \Sigma_1&=\Omega_0(W)(\gamma)
   =L(r\mapsto W^{[r+1]}(S(\gamma))),\\
 \Sigma_2&=\Omega_1(\Omega_0)(W)(\gamma).
\end{align*}
In the last expression, $\Omega_0$ is itself the function being iterated by
the next operator.  Expanded one step further, it is
\[
 \Sigma_2=L\bigl(r\mapsto
   [\Omega_0^{[r+1]}(x\mapsto S(W(x)))](\gamma)\bigr).
\]
Thus the second increase of level is not just another application of $W$.
It iterates the operation that constructs an orbit-limit transformation.
Higher levels continue this same typed pattern.  The recurrence for $D_k$
ensures that every such operator is eventually applied to appropriate
arguments and produces an ordinary member of $\mathcal T$.

\subsection{Feedback and the final number}
\label{amp:feedback}

Fix the evaluation argument $64$, and abbreviate
\[
 \sigma_b(a)=F_b(a;64),\qquad
 t_b(a)=\Sigma_{\sigma_b(a)+64}.
\]
The elementary advancement operation is
\begin{equation}
 A_b(a)=S(a)\star
 V\bigl(a\star t_b(a),\ t_b(a)\star a\bigr).
 \label{amp:eq:advance}
\end{equation}
Both orders of substitution occur in the payload.  The outer substitution
has first argument $S(a)$; that choice makes the successor gain available
before any properties of the payload are investigated.

Let $T$ be the tower of twos with the following exact height convention:
\[
 T(0)=1,\qquad T(r+1)=2^{T(r)}.
\]
One complete stage performs a number of advancement operations given by a
tower whose height is the current score plus one:
\begin{equation}
 G_b(a)=A_b^{[T(\sigma_b(a)+1)]}(a).
 \label{amp:eq:stage}
\end{equation}
Starting with $a_0=\Sigma_{64}$, define
\[
 a_{r+1}=G_b(a_r),\qquad s_b=\sigma_b(a_0).
\]
The proposed output of the amplifier is
\begin{equation}
 \boxed{N_b=\sigma_b(a_{s_b})=F_b(a_{s_b};64).}
 \label{amp:eq:candidate}
\end{equation}
The actual score of the initial tree is therefore also the number of
complete feedback stages.  At each stage, scores additionally control both
a tower height and the level of a function-space construction.

There is no circular definition here.  First define tree evaluation and
the tree operators.  Then define $A_b$ using the already available evaluator
on its input tree.  Finite iteration defines $G_b$, natural-number recursion
defines $a_r$, and the already defined initial score supplies the final
finite index.  During iteration, successive trees can become more complex;
termination of that iteration is nevertheless justified by its decreasing
natural-number counter.  A defining term need not reduce its enormous
numeric arguments during type checking.

There are two distinct uses of a score inside a stage.  Its total iteration
count $T(\sigma_b(a)+1)$ is fixed from the input tree before that stage
begins.  Within each advancement, the current tree determines a new
function-space index for $t_b$.  Recomputing that index does not reset the
iteration counter.  Likewise, the total number $s_b$ of complete stages is
fixed by the initial seed.  These distinctions rule out an interpretation
as an open-ended loop whose stopping condition continually moves away.
Every iteration count is a finite natural number supplied to an ordinary
recursor.

\subsection{A quantitative certificate independent of calibration}
\label{amp:bounds}

Assume for the remainder of this section that $b(n)\geq n+1$ everywhere.
Equation~\eqref{amp:eq:advance}, Proposition~\ref{amp:cor:graftlower}, and
Proposition~\ref{amp:prop:gain} give, for every evaluation argument $n$,
\begin{equation}
 F_b(A_b(a);n)\geq F_b(S(a);n)\geq F_b(a;n)+n.
 \label{amp:eq:advancegain}
\end{equation}
The payload in $A_b$ uses the score at $64$, but the displayed comparison
holds at every $n$.  Induction on the number of advancements now yields
\[
 F_b(A_b^{[r]}(a);n)\geq F_b(a;n)+rn.
\]
Writing $m=\sigma_b(a)$ and applying this inequality to one complete stage
gives the sharper intermediate estimate
\begin{equation}
 \sigma_b(G_b(a))\geq m+64T(m+1).
 \label{amp:eq:rawstage}
\end{equation}

The elementary inequalities $2^q\geq q+1$ imply by induction that
$T(r)\geq r+1$.  They also imply that $T$ is increasing.  Consequently
\eqref{amp:eq:rawstage} has two useful consequences:
\begin{align}
 \sigma_b(G_b(a))&\geq T(m+1),
 \label{amp:eq:stage-tower}\\
 \sigma_b(G_b(a))+1&\geq65(m+1).
 \label{amp:eq:stage-scale}
\end{align}
The first follows because $m\geq0$ and $64\geq1$.  For the second, use
$T(m+1)\geq m+1$ to obtain
$m+64T(m+1)+1\geq m+64(m+1)+1=65(m+1)$.

Define a purely numerical comparison function by
\[
 \Theta(s,0)=s,\qquad
 \Theta(s,r+1)=T(\Theta(s,r)+1).
\]
This repeatedly replaces the previous number by a tower of twos whose
height is that number plus one.  Its recursion is useful for communicating
a guaranteed lower bound: it requires none of the tree operators or higher
function spaces.

For example, the first two numerical feedback values are
\[
 \Theta(s,1)=T(s+1),\qquad
 \Theta(s,2)=T(T(s+1)+1).
\]
The second expression uses the entire first tower value as part of a new
tower height.  The final bound uses $s$ such replacements.  It is therefore
substantially more explicit than saying only that the candidate contains
many iterations: it gives a familiar recursively defined natural function
that the actual candidate is proved to exceed.

\begin{theorem}[Certified feedback growth]
\label{amp:thm:quantitative}
For every $r\in\mathbb N$,
\begin{align*}
 \sigma_b(a_r)&\geq\Theta(s_b,r),\\
 \sigma_b(a_r)+1&\geq65^r(s_b+1).
\end{align*}
In particular,
\[
 \boxed{N_b\geq\Theta(s_b,s_b)},\qquad
 \boxed{N_b+1\geq65^{s_b}(s_b+1)}.
\]
\end{theorem}

\begin{proof}
Both claims are proved by induction on $r$.  At $r=0$, the first is equality
by the definition of $s_b$, and the second uses $65^0=1$.  Suppose the first
claim holds at $r$.  Monotonicity of $T$, followed by
\eqref{amp:eq:stage-tower}, gives
\[
 \Theta(s_b,r+1)
 =T(\Theta(s_b,r)+1)
 \leq T(\sigma_b(a_r)+1)
 \leq\sigma_b(a_{r+1}).
\]
For the second claim, multiply its induction hypothesis by $65$ and apply
\eqref{amp:eq:stage-scale}.  Associativity of multiplication and the
successor equation for powers give the required factor $65^{r+1}$.
Finally specialize both inequalities to $r=s_b$.
\end{proof}

\begin{proposition}[Transferring a benchmark into the bound]
\label{amp:prop:transfer}
The function $\Theta$ is nondecreasing in each argument.  Hence, whenever
$c\leq s_b$,
\[
 \Theta(c,c)\leq N_b.
\]
Also $b(64)\leq s_b$ and $b(64)\leq N_b$.
\end{proposition}

\begin{proof}
Monotonicity in the first argument follows by induction on the second,
using monotonicity of $T$.  Since $T(q+1)\geq q+2$, each step of the
second argument increases the value, proving its monotonicity as well.
Thus $c\leq s_b$ gives
$\Theta(c,c)\leq\Theta(s_b,c)\leq\Theta(s_b,s_b)\leq N_b$.
The two base comparisons are instances of Lemma~\ref{amp:lem:base}, first
for $a_0$ and then for $a_{s_b}$.
\end{proof}

The proof distinguishes a guaranteed quantitative margin from a proposed
ordinal interpretation.  The entire payload appears in the concrete
definition of $N_b$, but the displayed lower bounds remain valid even if
its ordinal calibration is left unresolved.  They follow from structural
recursion, substitution, inflationarity, and explicit finite iteration.
This makes the benchmark comparison robust while leaving room for further
analysis of how much additional growth the higher operators contribute.

There is a useful modular consequence of the argument.  Suppose a later
investigation supplies a different explicit total tree transformation
$R_b(a)$.  Replacing the payload in~\eqref{amp:eq:advance} by $R_b(a)$ while
retaining the outer form $S(a)\star R_b(a)$ preserves
\eqref{amp:eq:advancegain}.  The induction giving the raw stage estimate and
both final bounds then applies without changing its reasoning.  The
substitution interface therefore isolates a stable mathematical obligation
from the choice of additional tree-generating operators.  Such a replacement
would still require its own well-typed definition and resource checks; the
present candidate is the fully specified payload above, rather than an
unspecified member of that wider family.

Conversely, the numerical bound alone does not identify the exact value of
$N_b$.  The raw stage inequality already discards most information about
the generated payload, and the tower comparison discards the additional
factor $64$ and the previous score.  These are deliberate weakenings in a
proof of a lower bound.  They make the certificate easy to state and
transport to a benchmark, while the exact recurrences retain all the
additional structure in the proposed number.


\section{The closed contender and its certified margin}
\label{sec:final}

Specialize the amplifier's base to the function \(b\) from
Section~\ref{sec:base}. Let \(s\) be its initial seed score and
\(\Cnew=\operatorname{candidate}(b)\). The complete tree construction is
fixed by the preceding definitions and reproduced as Lean code in the
appendix.

\begin{theorem}[Strictly beyond Contender 5]
\label{thm:main}
The supplied closed natural number satisfies
\[
 \Cfive<\Cnew.
\]
\end{theorem}
\begin{proof}
Proposition~\ref{prop:base} gives \(\Cfive<b(64)\). The amplifier's
universal base-domination lemma gives
\(b(64)\le F_b(a;64)\) for every tree \(a\), in particular for the final
feedback tree. Its value is \Cnew. Transitivity proves the claim.
\end{proof}

The corresponding checked proof is short because all enormous expressions
remain symbolic:
\begin{Verbatim}[frame=single,fontsize=\small,breaklines=true]
theorem contender_5_lt_contender_6 :
    Baseline.contender_5 < contender_6 :=
  Nat.lt_of_lt_of_le contender_5_lt_base64
    (Amplifier.base_le_candidate base base_weak)
\end{Verbatim}

\begin{theorem}[Incumbent-relative tower guarantee]
\label{thm:tower}
Define
\[
 T(0)=1,\qquad T(h+1)=2^{T(h)},
 \qquad
 \Theta(s,0)=s,\quad\Theta(s,k+1)=T(\Theta(s,k)+1).
\]
Then
\[
 \boxed{\Theta(\Cfive+1,\Cfive+1)\le\Cnew.}
\]
In addition,
\[
 65^{\Cfive+1}(\Cfive+2)\le\Cnew+1.
\]
\end{theorem}
\begin{proof}
The initial seed satisfies \(s\ge b(64)>\Cfive\), hence
\(\Cfive+1\le s\). The tower function is monotone; induction on the round
count shows that \(\Theta(s,k)\) is monotone in \(s\). It also increases
with the round count because \(T(x+1)\ge x+2\).
Thus
\[
 \Theta(\Cfive+1,\Cfive+1)\le\Theta(s,s)\le\Cnew,
\]
where the last inequality is the amplifier's feedback bound. For the second
claim, use its independently checked bound
\(65^s(s+1)\le\Cnew+1\), monotonicity of positive integer powers, and
\(\Cfive+1\le s\).
\end{proof}

Both claims are formalized, including the monotonicity needed to replace
\(s\) by a number expressed in terms of \Cfive. The power bound is a simple
readable consequence. The iterated-tower bound is the stronger headline
certificate. These are lower bounds, not asymptotic equivalents or a numerical
evaluation of the entry.

There is a practical proof-engineering detail here. Arithmetic lemmas should
first be proved for abstract natural-number variables and only then applied
to the enormous constants. In one direct closed instantiation, definitional
comparison attempted to unfold a large value through an arithmetic alias.
The final code uses a generic transfer lemma and, in the exponential bound,
spells the right-hand side as \code{Amplifier.candidate base + 1}. That is
definitionally \(\Cnew+1\), while avoiding unnecessary closed-value reduction.
No mathematical hypothesis or computational definition is changed by this
choice.

\section{What the tree notation does and does not assert}
\label{sec:ordinal}

The constructors named \code{omegaPower}, \code{veblenTree}, and
\code{gammaSeed} are explicit tree-building programs. Their names record
inspiration from familiar ordinal constructions. The comparison proofs use
their actual trees and the grafting identity, not a theorem identifying their
ordinal ranks. This distinction is essential for arbitrary raw trees.

\begin{proposition}[A counterexample to unrestricted ordinal calibration]
\label{prop:ordinalcounter}
Assign an ordinal value to a tree by
\[
 \rho(\mathsf{zero})=0,\quad
 \rho(\mathsf{succ}(a))=\rho(a)+1,\quad
 \rho(\mathsf{limit}(f))=\sup_{n<\omega}\rho(f(n)).
\]
The operation \code{veblenTree} does not, on all raw trees, factor through
this ordinal valuation as the standard Veblen function.
\end{proposition}
\begin{proof}
Let \(z=\mathsf{zero}\) and
\(\ell=\mathsf{limit}(n\mapsto z)\). Both have value zero. However,
\(\operatorname{veblenTree}(z,z)=\operatorname{omegaPower}(z)
=\mathsf{succ}(z)\), whose value is one.

For \(\ell\), the first-argument limit clause uses the common closure of
the constant family \(a\mapsto\operatorname{omegaPower}(a)\).
On input zero, its finite sweep of length \(r\) applies this operation
\(1+2+\cdots+r\) times. These lengths are unbounded. Structural induction
shows that \(\operatorname{omegaPower}\) has valuation
\(\alpha\mapsto\omega^\alpha\): grafting has valuation ordinal addition,
its finite repeated grafts give finite multiples, and the defining suprema
give the successor and limit laws for exponentiation. Hence the common
closure at zero has valuation
\[
 \sup_{j<\omega}\underbrace{\omega^{\omega^{\cdot^{\cdot^\omega}}}}
 _{\text{finitely many exponentiations}}=\varepsilon_0.
\]
One output has value one and the other has value \(\varepsilon_0\), although
their first inputs both have value zero and their second inputs agree.
Therefore the construction cannot generally depend only on the two ordinal
values in the asserted way.
\end{proof}

This is a mathematical analysis of a possible interpretation, not a Lean
ordinal formalization. In particular, the valuation \(\rho\) is not strict
well-founded tree height: a limit node may have the same value as its children.
The example does not affect totality or any numerical theorem in the package.
It identifies an assertion that would require additional restrictions and new
proofs before being made.

There is a second useful check: the formal theorem \code{fast\_at\_zero}
says \(F_b(a;0)=b(0)\) for \emph{every} tree \(a\). At a successor the
iteration count is one, and at a limit only branch zero is selected. Thus even
extremely complicated tree syntax supplies no strict gain at input zero.
Numerical dominance must follow from the evaluator's rules and hypotheses;
ordinal language alone cannot establish it. Our entry evaluates at 64, where
each explicit successor provides the proved positive gain.

\section{Validation and reproducibility}
\label{sec:validation}

\subsection{What was checked}

The archive contains a validator that stages the required Lean sources in a
fresh directory, discovers their local imports, and rebuilds every custom
dependency. It removes inherited project search paths, so pre-existing project
\code{.olean} files cannot silently supply the entry or proofs. Only the pinned
official toolchain's standard library is reused.

The definition phase contains \code{BaselineDefs}, \code{BaseDefs},
\code{AmplifierDefs}, and \code{CandidateDefs}. The proof phase contains
\code{Baseline}, \code{Base}, \code{Amplifier}, and \code{Candidate}.
An additional audit computes transitive axiom dependencies for 32 selected
public declarations, including all final comparison theorems. Every list is
empty. A declaration-body dependency traversal also verifies that
\code{contender\_6} does not depend on the closed fifth contender or either
of the source's selected-winner wrappers. It does use the general raw-language
interpreter and numerical envelope, as intended.

\begin{center}
\begin{tabular}{lrr}
\toprule
Phase & Measured seconds & Budget (seconds)\\
\midrule
All custom definitions & 1.408 & 15\\
All custom proofs & 1.363 & 60\\
Proofs plus axiom/dependency audits & 2.520 & 60\\
Additional stored-body kernel check & 2.458 & separate\\
\bottomrule
\end{tabular}
\end{center}
The required eight source modules contain 697 lines and 25,777 UTF-8 bytes. Compiler processes used one thread on an AMD EPYC 9V74 host running Ubuntu 24.04.3. The peak resident memory during definition/proof compilation was 428.3 MiB; including the extra checks, the peak was 927.9 MiB.
The additional check resubmitted 264 safe stored custom definition and theorem bodies directly to the C++ kernel under fresh names. It retained the imported standard-library declarations and inductive metadata; it is a body check, not a separate proof of the Lean implementation. Compiler-generated unsafe execution helpers were excluded, while the corresponding safe logical definitions were checked. The recorded scripts explain the precise scope.


These are observed wall times on the recorded environment, not hardware-
independent performance guarantees. The validator supports the same numerical
budgets on the reader's machine and fails rather than silently relaxing them.
No enormous natural number is evaluated during either budgeted phase.

The toolchain is Lean 4.19.0, commit
\begin{center}
\code{6caaee842e9495688c1567e78c0e68dbb96942aa}.
\end{center}
The official compiler and kernel libraries were unchanged. This container
required a small disclosed application-path launcher because its ordinary
executable-path query was unavailable. The launcher supplies its own resolved
program path; it does not modify the elaborator or kernel. Its complete C
source, build instructions, release archive hash, and shared-library hash are
included in \artifact{scripts/REPRODUCTION.md} and the adjacent source files.
An ordinary official Lean installation does not need this launcher.

\subsection{Independent semantic checks}

The Python model in \artifact{scripts/reference_model.py} independently
implements the original raw syntax, cast pairing, reversed lookup, and
recursor. Its depth-zero-through-four rows agree with Lean's executable smoke
tests. Both implementations also check an example of recursion returning a
function: start with successor, double function composition three times, and
apply the resulting function to zero; the answer is eight. The Lean version
is a closed reduction proof. It uses a locally increased reduction-depth
limit only in this optional test, not in the entry or comparison modules.

The independent model is a translation check, not a second proof assistant.
Likewise, finite tests cannot prove that two implementations agree at all
inputs. The universal comparison theorem is supplied by Lean, and the
cross-language correspondence is supported by the detailed source audit and
literal preservation of the relevant computational definitions.

\subsection{Commands}

With the pinned Lean toolchain available, run from the archive root:
\begin{Verbatim}[frame=single,fontsize=\small,breaklines=true]
python3 scripts/validate.py --lean lean
python3 scripts/reference_model.py
lake build
\end{Verbatim}
The first command is the budgeted, clean-source validation. The third is a
convenient ordinary project build; it is not the source of the timing claim.
The optional executable smoke modules are deliberately outside the submitted
entry's import graph.

To rebuild this article, run
\begin{Verbatim}[frame=single,fontsize=\small]
cd article
latexmk -xelatex -interaction=nonstopmode -halt-on-error \
  Lean_Big_Number.tex
\end{Verbatim}
The source uses standard TeX packages, Latin Modern fonts, and DejaVu Sans Mono
for Unicode Lean listings. All input fragments are supplied.

\section{Constructivity, the diagonal barrier, and optimality}

\subsection{Why the definition is a natural number}

Every component is a total definition in the constructive fragment used by the
submission. The language enumerators recurse on natural depth. Interpretation
recurses on syntax and uses ordinary natural-number recursors at interpreted
types. The tree evaluator recurses on the inductive tree; its extra iterations
have finite natural counts. Higher function spaces are formed by finite
recursion on a natural level. Feedback uses a value from a previously
constructed tree to choose a new finite count or level. There is no equation
that defines a tree or natural number by presupposing its own final value.

A limit constructor takes a function \(\N\to\mathsf{Tree}\). This does not
require writing infinitely many subtrees into a file. A finite, total program
provides the children. Structural recursion receives each requested child's
recursive result, as licensed by the inductive recursor. A cyclic definition
would not pass the same termination check.

The finite maximum defining \(E(n)\) ranges over a language with an explicitly
given total evaluator. It does not search arbitrary programs and wait for an
unknown set of halting computations. Accordingly, the entry does not depend
on a Busy Beaver bound, a halting oracle, classical selection of an unspecified
maximum, or an unproved consistency assumption.

\subsection{Why one cannot simply maximize over all Lean definitions}

A tempting shortcut is to replace the object-language envelope by a purported
total evaluator for every total Lean function, then diagonalize. That step
requires a justification which the present construction deliberately does
not assume. The elementary obstruction is worth stating precisely.

\begin{proposition}[Diagonal obstruction]
Suppose an enumerated class of total unary functions has a total evaluator
\(U(e,n)\), represents every function in that class by some code, and is closed
under the function \(d(n)=U(n,n)+1\) defined in the same class. These assumptions
are inconsistent.
\end{proposition}
\begin{proof}
If \(e\) represents \(d\), then
\(d(e)=U(e,e)+1=d(e)+1\), impossible in \(\N\).
\end{proof}

This proposition does not rule out total interpretation of a weaker object
language inside a stronger host language. That is exactly the relationship
used here: the raw system-T-style language is interpreted in Lean, and the
new tree feedback is an additional host-language construction. The argument
also does not forbid fixed, finite maxima with independently justified total
interpreters. It identifies why a supposed universal same-language evaluator
cannot be inserted without resolving its scope and totality obligations.

\subsection{The meaning of ``largest''}

This article proposes a concrete entry and proves its performance against the
specified incumbent. It does not establish that every admissible Lean program
names a smaller number. If a fixed machine, time bound, and finite source-size
cap produce a finite set of accepted source strings, then that finite contest
has a maximum score abstractly. The work here neither characterizes that set
nor proves optimality over it.

The elementary successor operation explains why there is no largest natural
number, and can improve a particular submission when sufficient resource
slack remains. It is not, by itself, a proof that a resource-bounded contest
has no maximum. Nor would it justify presenting a circular ``one greater than
every permissible entry'' description as a constructive Lean definition.
The present submission has an explicit source, a total computation, and a
checked strict margin; those are the claimed achievements.

\section{Further research directions}

\begin{enumerate}
\item \textbf{A calibrated domain of ordinal trees.} Identify a decidable or
constructively checkable class of well-behaved trees on which the closure
constructors respect an ordinal valuation. The degenerate-limit counterexample
shows that restrictions must be stated, not inferred from notation. A useful
result would combine closure of the class under the actual feedback operators
with a proved numerical comparison principle.
\item \textbf{Direct majorization of the raw interpreter.} Prove bounds for
\(E(n)\) in a separately specified fast-growing hierarchy, including the effects
of type annotations, recursor result types, malformed terms, and default
casts. Such a theorem could replace the exhaustive envelope in the base by a
more transparent closed growth function while retaining a checked comparison
with the exact fifth contender.
\item \textbf{A richer certified object language.} Extend the finite syntax
with one explicit inductive tree type and its recursor, give a total
interpretation, and prove an embedding of the current raw language with a
controlled depth overhead. Iterating carefully stratified language extensions
could provide a different route to stronger contenders without positing a
universal Lean self-interpreter.
\item \textbf{A verified translation bridge.} Formalize enough of the Coq
computational subset to certify the correspondence of the two implementations.
An intermediate first-order syntax for types, raw terms, casts, and evaluator
equations would make the remaining cross-language trust boundary explicit.
\item \textbf{Sharper structural lower bounds.} The graft proof intentionally
ignores most of the Veblen-inspired payload. Find syntax-sensitive order
relations preserved by evaluation at controlled inputs, and use them to
quantify the higher-type seeds. The equality at input zero warns against
unqualified pointwise monotonicity claims.
\item \textbf{Checking efficiency with strict dependency audits.} Explore
proof organization, abstract interfaces, and kernel conversion behavior while
preserving empty axiom lists. The closed-value alias issue illustrates why
symbolic theorem design can matter more than the physical length of a proof.
Record total custom dependency cost, not just the final declaration's cost.
\end{enumerate}

\appendix
\section{Archive map and provenance}

\begin{longtable}{@{}p{0.38\textwidth}p{0.57\textwidth}@{}}
\toprule
Path & Purpose\\
\midrule
\endhead
\artifact{BigNumber/BaselineDefs.lean} & Literal computational port and numerical envelope.\\
\artifact{BigNumber/Baseline.lean} & Selected-max invariant and axiom-free maximum proofs.\\
\artifact{BigNumber/BaseDefs.lean} & Inclusive envelope sum and independent base.\\
\artifact{BigNumber/Base.lean} & Base domination, including the exact incumbent.\\
\artifact{BigNumber/AmplifierDefs.lean} & Trees, interpreter, closure constructors, higher types, feedback.\\
\artifact{BigNumber/Amplifier.lean} & Grafting, gains, monotonicity, and quantitative lower bounds.\\
\artifact{BigNumber/CandidateDefs.lean} & The submitted closed natural number.\\
\artifact{BigNumber/Candidate.lean} & Final strict and incumbent-relative certificates.\\
\artifact{BigNumber/BaselineSmoke.lean} & Optional executable small-depth tests.\\
\artifact{BigNumber/SemanticChecks.lean} & Closed higher-type recursor reduction test.\\
\artifact{scripts/validate.py} & Clean build, cumulative budgets, axiom and dependency audit.\\
\artifact{scripts/reference_model.py} & Independent small-input semantic model.\\
\artifact{validation/} & Recorded outcomes, logs, and independent mathematical review.\\
\artifact{upstream/} & Pinned Coq source, original rules, and MIT license.\\
\artifact{article/} & Complete article source, fragments, and compiled PDF.\\
\bottomrule
\end{longtable}

The upstream \code{Contender.v} Git blob is
\code{f92db72deb0f3a4d4e3b575b4949eda73b1417a0}. The exact source files and their
SHA-256 hashes are recorded in the archive manifest. The original source is
redistributed under its MIT license, retained in \artifact{upstream/LICENSE}.
New material is separately identified in the root README and license.

\section{Complete computational amplifier}

This listing is the exact computational module used in validation. It has no
custom imports; the standard Lean prelude is implicit. The baseline and base
modules are supplied as separate source files.
\VerbatimInput[fontsize=\scriptsize,breaklines=true,breakanywhere=true,
 numbers=left,numbersep=5pt]{../BigNumber/AmplifierDefs.lean}

\section{The base and final comparison proofs}

The following are the exact base and final proof modules. All supporting
proofs, including the full graft and feedback theorems, are supplied in the
archive and explained in the main text.
\subsection{Base definitions}
\VerbatimInput[fontsize=\scriptsize,breaklines=true,breakanywhere=true,
 numbers=left,numbersep=5pt]{../BigNumber/BaseDefs.lean}
\subsection{Base proofs}
\VerbatimInput[fontsize=\scriptsize,breaklines=true,breakanywhere=true,
 numbers=left,numbersep=5pt]{../BigNumber/Base.lean}
\subsection{Closed definition}
\VerbatimInput[fontsize=\scriptsize,breaklines=true,breakanywhere=true,
 numbers=left,numbersep=5pt]{../BigNumber/CandidateDefs.lean}
\subsection{Final certificates}
\VerbatimInput[fontsize=\scriptsize,breaklines=true,breakanywhere=true,
 numbers=left,numbersep=5pt]{../BigNumber/Candidate.lean}

\begin{thebibliography}{9}
\raggedright
\bibitem{rules}
Cody Roux and contributors.
\emph{name-the-biggest-number: contest rules}.
Pinned repository snapshot, accessed 12 October 2026.
\url{https://github.com/codyroux/name-the-biggest-number/blob/d56dce013e6264292533efdc5e1de71863c6fb86/README.md}.

\bibitem{contender}
Cody Roux, Samuel Gruetter, and contributors.
\emph{Contender.v}, especially the raw-language evaluator and Contender 5.
Pinned repository snapshot.
\url{https://github.com/codyroux/name-the-biggest-number/blob/d56dce013e6264292533efdc5e1de71863c6fb86/Contender.v}.

\bibitem{leanaxioms}
Lean community.
\emph{Theorem Proving in Lean 4}, chapter ``Axioms and Computation.''
Accessed 12 October 2026.
\url{https://lean-lang.org/theorem_proving_in_lean4/Axioms-and-Computation/}.

\bibitem{leanrelease}
Lean development team.
\emph{Lean 4.19.0}, pinned release and implementation.
\url{https://github.com/leanprover/lean4/releases/tag/v4.19.0}.
The exact implementation commit and library hashes used in this work appear
in the reproduction record.
\end{thebibliography}
\end{document}
